\documentclass[11pt]{article}
\usepackage[a4paper,top=1.4cm,bottom=1.3cm,left=1.7cm,right=1.7cm]{geometry}
\usepackage{fontspec}
\usepackage{polyglossia}
\setmainlanguage{russian}
\setotherlanguage{english}
\setmainfont{DejaVu Serif}[Scale=0.85]
\setsansfont{DejaVu Sans}[Scale=0.9]
\setmonofont{DejaVu Sans Mono}[Scale=0.76]
\usepackage{setspace}\setstretch{0.90}
\setlength{\abovedisplayskip}{5pt}\setlength{\belowdisplayskip}{5pt}
\usepackage{amsmath,amssymb}
\usepackage{booktabs}
\usepackage{enumitem}
\setlist{nosep,leftmargin=1.4em}

\title{\vspace{-1.4cm}Методика: матрицы чувствительности 19$\times$19\\[2pt]
\large Замена входных полей на среднеклиматологические}
\author{}
\date{}

\begin{document}
\maketitle
\vspace{-1.2cm}

\section{Формат данных}

\paragraph{Сетка.} Все поля заданы на регулярной широтно-долготной сетке $0.25^\circ$:
$N_\lambda = 1440$ долгот, $N_\varphi = 721$ широта. Aurora обрезает южный полюс, поэтому
её выходы и сравниваемые с ними поля берутся на $N_\varphi = 720$; Pangu работает на полной
сетке 721. Обозначим множество узлов $\Omega$, $N = |\Omega| = N_\varphi N_\lambda$.

\paragraph{Переменные.} Рассматриваются $|\mathcal{V}| = 19$ полей~--- ровно те, для которых
есть скачанная климатология, поэтому подмена всегда полная, без частичных столбцов:
\begin{itemize}
\item приземные (4): \texttt{MSLP}, \texttt{U10}, \texttt{V10}, \texttt{T2M};
\item высотные (15): $Z$, $Q$, $T$, $U$, $V$ на уровнях 1000, 850, 500 гПа.
\end{itemize}
Обозначение $Z$ здесь~--- \emph{геопотенциал} $\Phi$ (переменная ERA5 \texttt{geopotential},
короткое имя \texttt{z}), единицы м$^2$/с$^2$. Это не геопотенциальная высота
$Z_g = \Phi/g_0$ (гпм) при $g_0 = 9.80665$~м/с$^2$: деление на $g_0$ не выполняется. На значения
в матрицах это не влияет, поскольку обе метрики инвариантны к постоянному множителю~---
$S_{vu}$ нормирована на СКО того же поля, а $\Delta\!\operatorname{ACC}$ безразмерна.

\paragraph{Взвешивание по площади.} Узлы регулярной широтно-долготной сетки покрывают
существенно разную площадь: ячейка на широте $89.75^\circ$ меньше экваториальной примерно
в $\cos(89.75^\circ)\approx 0.004$ раза, но в невзвешенной сумме весила бы столько же. Поэтому
все суммы ниже берутся с весом $w_s = \cos\varphi_s$ (нормированным на единицу) --- принятая
в WeatherBench-2 и ECMWF конвенция:
\[
w_s = \frac{\cos\varphi_s}{\sum_{s'\in\Omega}\cos\varphi_{s'}},
\qquad \sum_{s\in\Omega} w_s = 1 .
\]

\paragraph{Истина.} ERA5, 6-часовой шаг: $y_u(t) \in \mathbb{R}^{\Omega}$~--- поле переменной
$u$ в момент $t$.

\paragraph{Климатология.} WeatherBench-2, предпосчитанная по 1990--2019 как среднее по годам
для каждой пары (день года, час суток) с последующим сглаживанием:
\[
\bar c_u(t) \;\approx\; \frac{1}{30}\sum_{y=1990}^{2019} y_u\big(y,\, d(t),\, h(t)\big),
\qquad d(t)\in[1,366],\quad h(t)\in\{0,6,12,18\}.
\]
Ключевая деталь: климатология зависит и от дня года, и от часа~--- то есть содержит суточный ход.

\paragraph{Инициализации.} $\mathcal{D}$~--- 48 дат за 2024 год: 4-е, 11-е, 18-е и 25-е числа
каждого из 12 месяцев, час инициализации циклически 00/06/12/18 (по 12 раз каждый), чтобы
сезон и время суток были представлены равномерно.

\section{Постановка эксперимента}

Пусть $\mathbf{x}(t) = \{x_u(t)\}_{u}$~--- полное входное состояние модели. Aurora принимает
\emph{два} снимка $\big(\mathbf{x}(t_0-6\text{ч}),\,\mathbf{x}(t_0)\big)$, Pangu~--- \emph{один}
$\mathbf{x}(t_0)$.

\paragraph{Оператор подмены.} Для переменной $v \in \mathcal{V}$ определим
\[
\big(P_v\mathbf{x}\big)_u(t) \;=\;
\begin{cases}
\bar c_v(t), & u = v,\\[2pt]
x_u(t), & u \neq v,
\end{cases}
\qquad \text{глобально, во всех узлах } \Omega .
\]
Для Aurora подмена применяется к обоим входным снимкам, для Pangu~--- к единственному.

\paragraph{Прогноз.} Пусть $\mathcal{M}$~--- один шаг модели на 6 часов, применяемый
авторегрессивно. Для лида $\tau \in \{6, 24\}$ часов:
\[
\hat y^{(0)}_u(t_0,\tau) \;=\; \big[\mathcal{M}^{\tau/6}\mathbf{x}(t_0)\big]_u
\quad\text{(базовый прогон)},
\qquad
\hat y^{(v)}_u(t_0,\tau) \;=\; \big[\mathcal{M}^{\tau/6}P_v\mathbf{x}(t_0)\big]_u .
\]
На каждую дату выполняется $1 + 19 = 20$ прогонов; из каждого сохраняются все 19 полей на
двух лидах, то есть $20 \times 2 \times 19 = 760$ массивов.

\section{Метрики на графиках}

Обе матрицы имеют размер $19\times19$: \textbf{строка} $v$~--- какую переменную подменили,
\textbf{столбец} $u$~--- на какую переменную смотрим в прогнозе. Диагональ $v = u$~---
самовлияние.

\subsection*{Верхний ряд: относительная чувствительность}

Сравнивается патченый прогон с базовым (истина используется только для нормировки):
\[
\boxed{\;
S_{vu}(\tau) \;=\; \frac{1}{|\mathcal{D}|}\sum_{t_0\in\mathcal{D}}
\frac{\operatorname{RMSE}_w\!\big(\hat y^{(v)}_u(t_0,\tau),\; \hat y^{(0)}_u(t_0,\tau)\big)}
     {\sigma_w\big[y_u(t_0+\tau)\big]}
\;}
\]
где
\[
\operatorname{RMSE}_w(a,b) = \sqrt{\sum_{s\in\Omega} w_s\big(a_s-b_s\big)^2},
\qquad
\sigma_w[a] = \sqrt{\sum_{s\in\Omega} w_s\big(a_s-\bar a_w\big)^2},\quad
\bar a_w = \sum_{s\in\Omega} w_s\, a_s .
\]
Деление на пространственное СКО \emph{истинного} поля делает величину безразмерной, поэтому
клетки с разными единицами (Па, м$^2$/с$^2$, кг/кг, К, м/с) сопоставимы. Смысл: во сколько раз возмущение
прогноза от подмены $v$ велико по сравнению с естественной изменчивостью поля $u$.
Усреднение по 48 датам выполняется \emph{после} нормировки, на уровне отдельной даты.

\subsection*{Нижний ряд: изменение ACC}

Аномальная корреляция (anomaly correlation coefficient) прогноза $\hat y$ с истиной $y$
относительно климатологии $\bar c$:
\[
\operatorname{ACC}_w(\hat y, y, \bar c) \;=\;
\frac{\sum_{s\in\Omega} w_s\big(\hat y_s - \bar c_s\big)\big(y_s - \bar c_s\big)}
     {\sqrt{\sum_{s\in\Omega} w_s\big(\hat y_s - \bar c_s\big)^2}\;
      \sqrt{\sum_{s\in\Omega} w_s\big(y_s - \bar c_s\big)^2}} \;\in[-1,1].
\]
На графике показано, насколько подмена ухудшает эту величину относительно базового прогона:
\[
\boxed{\;
\Delta\!\operatorname{ACC}_{vu}(\tau) \;=\; \frac{1}{|\mathcal{D}|}\sum_{t_0\in\mathcal{D}}
\Big[\operatorname{ACC}_w\big(\hat y^{(v)}_u,\, y_u,\, \bar c_u\big)
   - \operatorname{ACC}_w\big(\hat y^{(0)}_u,\, y_u,\, \bar c_u\big)\Big]
\;}
\]
Отрицательные значения означают потерю качества. В отличие от $S_{vu}$, эта метрика опирается
на истину и отвечает на вопрос «стало ли \emph{хуже}», а не просто «изменилось ли».

\subsection*{Что важно помнить при чтении}

\begin{itemize}
\item Обе метрики взвешены по площади ячейки ($w_s\propto\cos\varphi_s$). Проверка показала,
      что переход от невзвешенной версии к взвешенной практически не меняет структуру матриц:
      корреляция внедиагональных элементов $r \approx 0.99$ для обеих моделей, суммарная
      внедиагональная «масса» растёт на 1--4\,\%. То есть исправление нужно для соответствия
      стандарту, а не потому, что оно меняет выводы.
\item В $\operatorname{ACC}_w$ среднее по полю не вычитается~--- это стандартное определение,
      где роль «среднего» играет климатология $\bar c$.
\item $S_{vu}$ нормирована на $\sigma_w$ \emph{столбца} $u$. Сравнение столбцов между собой
      поэтому смещено в пользу малоизменчивых полей (приземный ветер); сравнение пары
      $v\!\to\!u$ против $u\!\to\!v$ от этого свободно, так как нормировки различаются.
\item Это чувствительность к \emph{конечному} возмущению (поле целиком заменяется
      климатологией), далеко за пределами линейного режима, а не локальный градиент.
\item Возмущения у моделей не равны: Aurora получает два входных снимка (подменяются оба),
      Pangu~--- один. Это следствие архитектур, а не методики.
\end{itemize}

\end{document}
